\chapter{Vocabulario}

\newcounter{phrase}
\newcounter{beginning}
\newcounter{ending}
\newcounter{sequence}

%pagina 715 de Arab Music Theory in the modern period

La música árabe, y en especial en la improvisación, hace uso de frases típicas que forman parte de la tradición.

En este capítulo veremos una colección de frases que pueden considerarse <<coloquiales>>, y cuyo conocimiento es es una herramienta fundamental para lograr que una improvisación suene natural y genuina, dentro de lo que dicta la tradición.

Por sí solo, el contenido de este capítulo no sirve para elaborar un vocabulario completo y, más importante aún, saber cómo y cuando emplearlo, de la misma manera que un libro que recoja la jerga de un idioma no basta para que un hablante sepa cómo utilizar esta. Es necesaria la exposición a esa jerga en su contexto real, verla puesta en práctica, para inferir a partir de ello su utilización, de la misma manera que es necesario escuchar las frases de este capítulo dentro de sus correspondientes improvisaciones para después ser capaz de incorporarlas en las propias de manera natural. No obstante, una recopilación así de patrones melódicos a la manera de un diccionario resulta útil para disponer de un vocabulario básico y poder de este modo entender mejor esas improvisaciones ajenas, haciendo más productivo el estudio de \taqasim y piezas de otros intérpretes.

No deben entenderse estos fragmentos como frases que han de tocarse literalmente y sin ningún cambio respecto a lo que aquí se presenta. Al contrario, deben servir como base para que cada cual elabore sus propios motivos melódicos y dé así forma a su estilo. Y, por supuesto, incluso en el caso de tocar la melodía de la frase tal y como se recoge, se puede (y se debe) añadir la ornamentación correspondiente.

Dividiremos los patrones de este capítulo en tres grupos principales: inicios, cierres, frases breves y secuencias melódicas.

Los ejemplos comentados a continuación se presentan en su  mayoría en los \maqamaat Rast y Nahawand, salvo cuando se indica lo contrario. El objetivo no es solo ofrecer una colección de fragmentos y frases, sino presentar en base a ellos las formas más habituales en las que se presentan como elementos de una pieza, para así ayudar a una mejor comprensión de la estructura de un \taqsim y de otras frases similares. Junto a ellos, y en las dos primeras secciones (inicios y cierres), se incluyen secciones particulares para algunos de los \maqamaat más habituales.

\section{Inicios}

En esta sección se recogen fragmentos con los que es habitual comenzar un \taqsimb, y que siven de presentación de su \maqamb. La apertura y cierre de un \taqsim son las partes que más se prestan al uso de <<formulas>> predefinidas, y donde el conocimiento de estas es más necesario.

La manera más prototípica de empezar un \taqsim o una frase dentro de este es simplemente repetir la nota tónica u otra de sus notas a la que se quiera dar relevancia. Habitualmente, se utiliza también la nota tónica en una octava inferior, intercalandola entre las repeticiones de la anterior.


\begin{center}
\begin{lilypond}
\include "arabic.ly"
notes = \relative {
\key do \rast
\time 16/4
do8\downbow do do do do do do do do, \downbow do'\downbow do do do do do do
\bar "||"
}

\layout {
  \omit Voice.StringNumber
  \context {
      \TabStaff
      \tabFullNotation
      stringTunings = \stringTuning <do, fa, la, re sol do'>
    }
  indent = #0
  \context {
    \Score
    defaultBarType = ""
    supportNonIntegerFret = ##t
  }
  \context {
    \TabStaff
    \remove Time_signature_engraver

    }
}

<<
  \new TabStaff <<\notes>>
>>
\end{lilypond}
\end{center}


\audioclip{Inicio con nota tónica.}

Nótese que las distintas pulsaciones sobre la misma cuerda se hacen con una pulsación alterna, pero tanto la pulsación de la cuerda más grave como la de la nota sucesiva una octava por encima se han de hacer con pulsación descendente. Se aplica aquí el principio de que, siempre que se produzca un cambio de cuerda, la primera pulsación sobre esta nueva cuerda ha de ser en sentido descendente.

En lugar de comenzar directamente con la nota tónica, es habitual llegar hasta ella tras una pequeña frase que presente el \yins en cuestión. El siguiente inicio es prototípico en el \maqam Bayati:


\begin{center}
\begin{lilypond}
\include "arabic.ly"
notes = \relative {
\key re \bayati
fa8\downbow misb re do re re re re re re re re re,\downbow re'\downbow re re
\bar "||"
}

\layout {
  \omit Voice.StringNumber
  \context {
      \TabStaff
      \tabFullNotation
      stringTunings = \stringTuning <do, fa, la, re sol do'>
    }
  indent = #0
  \context {
    \Score
    defaultBarType = ""
    supportNonIntegerFret = ##t
  }
  \context {
    \TabStaff
    \remove Time_signature_engraver

    }
}

<<
  \new TabStaff <<\notes>>
>>
\end{lilypond}
\end{center}


\audioclip{Inicio con frase antes de llegar a tónica (Bayati).}

El siguiente es un ejemplo similar en \maqam Rast, que en este caso no llega a la tónica sino a la cuarta nota del \maqam (fa). Puede emplearse para arrancar una frase intermedia de un \taqsimb, en la que quiera enfatizarse esta nota.


\begin{center}
\begin{lilypond}
\include "arabic.ly"
notes = \relative {
\key do \rast
  do,8 do' re misb misb fa misb re misb fa fa fa fa fa fa fa
\bar "||"
}

\layout {
  \omit Voice.StringNumber
  \context {
      \TabStaff
      \tabFullNotation
      stringTunings = \stringTuning <do, fa, la, re sol do'>
    }
  indent = #0
  \context {
    \Score
    defaultBarType = ""
    supportNonIntegerFret = ##t
  }
  \context {
    \TabStaff
    \remove Time_signature_engraver

    }
}

<<
  \new TabStaff <<\notes>>
>>
\end{lilypond}
\end{center}


\audioclip{Inicio con frase.}

Es frecuente que la frase de apertura utilice las notas por debajo de la tónica, como en el siguiente ejemplo:


\begin{center}
\begin{lilypond}
\include "arabic.ly"
notes = \relative {
\key do \rast
  r8 sol,8 la sisb do do do do do, do' do do
\bar "||"
}

\layout {
  \omit Voice.StringNumber
  \context {
      \TabStaff
      \tabFullNotation
      stringTunings = \stringTuning <do, fa, la, re sol do'>
    }
  indent = #0
  \context {
    \Score
    defaultBarType = ""
    supportNonIntegerFret = ##t
  }
  \context {
    \TabStaff
    \remove Time_signature_engraver

    }
}

<<
  \new TabStaff <<\notes>>
>>
\end{lilypond}
\end{center}


\audioclip{Inicio llegando a tónica desde octava inferior.}

En el caso de \maqam Bayati, a pesar de que su segundo \yins es \yins Nahawand en sol, se suele utilizar en la octava inferior \yins Rast en sol (es decir, con si\hflat\ en lugar de si$\flat$). El siguiente es un ejemplo de arranque en \maqam Bayati utilizando notas por debajo de la tónica.


\begin{center}
\begin{lilypond}
\include "arabic.ly"
notes = \relative {
\key re \bayati
sol8 fa misb re do sisb do re re, re' re re, re' re re, re'
\bar "||"
}

\layout {
  \omit Voice.StringNumber
  \context {
      \TabStaff
      \tabFullNotation
      stringTunings = \stringTuning <do, fa, la, re sol do'>
    }
  indent = #0
  \context {
    \Score
    defaultBarType = ""
    supportNonIntegerFret = ##t
  }
  \context {
    \TabStaff
    \remove Time_signature_engraver

    }
}

<<
  \new TabStaff <<\notes>>
>>
\end{lilypond}
\end{center}


\audioclip{Inicio llegando a tónica desde octava inferior (Bayati).}

La primera frase de un \taqsim suele ser breve y presentar el primer \yins del \maqam, de tal modo que muchas veces no es apenas más que la combinación de un inicio y un cierre. Veáse el ejemplo siguiente:


\begin{center}
\begin{lilypond}
\include "arabic.ly"
notes = \relative {
\key do \rast
do,8 do' do do do do do do sisb32 do sisb la sisb8 sisb do do re misb16 re do8
\bar "|"
\break
re misb fa fa misb misb re re
misb16 re
do2
\bar "||"
}

\layout {
  \omit Voice.StringNumber
  \context {
      \TabStaff
      \tabFullNotation
      stringTunings = \stringTuning <do, fa, la, re sol do'>
    }
  indent = #0
  \context {
    \Score
    defaultBarType = ""
    supportNonIntegerFret = ##t
  }
  \context {
    \TabStaff
    \remove Time_signature_engraver

    }
}

<<
  \new TabStaff <<\notes>>
>>
\end{lilypond}
\end{center}


\audioclip{Primera frase completa.}

O este otro en \maqam Nahawand, más breve, que aun así sirve para inicar con claridad que el \taqsim se desarrolla en este \maqam:


\begin{center}
\begin{lilypond}
\include "arabic.ly"
notes = \relative {
mib16 fa sol8 fa sol fa16 mib \tuplet 3/2 {mib (fa mib)} \tuplet 3/2 {re (mib re)} do4
\bar "||"
}

\layout {
  \omit Voice.StringNumber
  \context {
      \TabStaff
      \tabFullNotation
      stringTunings = \stringTuning <do, fa, la, re sol do'>
    }
  indent = #0
  \context {
    \Score
    defaultBarType = ""
    supportNonIntegerFret = ##t
  }
  \context {
    \TabStaff
    \remove Time_signature_engraver

    }
}

<<
  \new TabStaff <<\notes>>
>>
\end{lilypond}
\end{center}


\audioclip{Primera frase completa.}

Una técnica muy habitual en un \taqsim es tocar una frase breve como las anteriores, formada por un inicio y un cierre, e iniciar la frase siguiente, mas larga, repitiendo el mismo inicio. Esto puede verse en el siguiente ejemplo en \maqam Bayati.


\begin{center}
\begin{lilypond}
\include "arabic.ly"
notes = \relative {
\key re \bayati
\partial 16 do,16
re re' re re re misb32 (re) do16 re
re,16 re' re re re misb32 (re) do16 re
re,16 re' re re re r2
\bar "|"
\break
do,16
re re' re re re misb32 (re) do16 re
re,16 re' re re re misb32 (re) do16 re
re, re' misb misb fa fa misb misb8
fa8 fa8:32 misb16 misb (re) do re re4
\bar "||"
}

\layout {
  \omit Voice.StringNumber
  \context {
      \TabStaff
      \tabFullNotation
      stringTunings = \stringTuning <do, fa, la, re sol do'>
    }
  indent = #0
  \context {
    \Score
    defaultBarType = ""
    supportNonIntegerFret = ##t
  }
  \context {
    \TabStaff
    \remove Time_signature_engraver

    }
}

<<
  \new TabStaff <<\notes>>
>>
\end{lilypond}
\end{center}


\audioclip{Primera frase completa y segunda frase con mismo inicio.}

En ocasiones, esa primera frase no contiene un cierre, sino que se deja inconclusa sobre otra nota del \maqam, creándose un efecto de tensión que se resuelve cuando la segunda frase arranca. Por ejemplo, en el siguiente ejemplo en \maqam Rast, la primera frase se detiene sobre la nota fa.


\begin{center}
\begin{lilypond}
\include "arabic.ly"
notes = \relative {
\time 2/4
do,16 do' do do do do do do
do,16 do' do do do sisb32 sisb sisb16 do32 do do16 re16 re misb misb fa
\override TextSpanner.bound-details.left.text = "rit."
\startTextSpan
fa fa \bar "|" fa fa fa fa
\stopTextSpan
r4
do,16 do' do do do do do do
do,16 do' do do do sisb32 sisb sisb16 do32 do \bar "|"
do16 re16 re misb misb fa fa
misb misb re re16 misb32 re do16 re re misb misb fa fa sol sol fa^"etc." fa misb
\bar "||"
}

\layout {
  \omit Voice.StringNumber
  \context {
      \TabStaff
      \tabFullNotation
      stringTunings = \stringTuning <do, fa, la, re sol do'>
    }
  indent = #0
  \context {
    \Score
    defaultBarType = ""
    supportNonIntegerFret = ##t
  }
  \context {
    \TabStaff
    \remove Time_signature_engraver

    }
}

<<
  \new TabStaff <<\notes>>
>>
\end{lilypond}
\end{center}


\audioclip{Primera frase con final inconcluso.}

Los inicios anteriores concluyen todos ellos sobre la tónica. Sin embargo, una frase de inicio puede servir para llegar hasta cualquier otra nota del \maqamb. Un caso particular son los que arrancan de la tónica y ascienden hasta el \emph{ghammaz}, para allí dar inicio a una frase que se situa sobre el segundo \yins. Estos inicios no es habitual utilizarlos para comenzar el \taqsimb, sino como arranque de frases intermedias.

Los siguientes son algunos ejemplos.


\begin{center}
\begin{lilypond}
\include "arabic.ly"
notes = \relative {
\key do \rast
do16 sisb do re misb re misb fa
sol8 sol sol sol
\bar "||"
}

\layout {
  \omit Voice.StringNumber
  \context {
      \TabStaff
      \tabFullNotation
      stringTunings = \stringTuning <do, fa, la, re sol do'>
    }
  indent = #0
  \context {
    \Score
    defaultBarType = ""
    supportNonIntegerFret = ##t
  }
  \context {
    \TabStaff
    \remove Time_signature_engraver

    }
}

<<
  \new TabStaff <<\notes>>
>>
\end{lilypond}
\end{center}


\audioclip{Inicio ascendente hasta \emph{ghammaz}.}


\begin{center}
\begin{lilypond}
\include "arabic.ly"
notes = \relative {
\key do \rast
do16 do re re misb misb fa fa
sol8 sol sol sol
\bar "||"
}

\layout {
  \omit Voice.StringNumber
  \context {
      \TabStaff
      \tabFullNotation
      stringTunings = \stringTuning <do, fa, la, re sol do'>
    }
  indent = #0
  \context {
    \Score
    defaultBarType = ""
    supportNonIntegerFret = ##t
  }
  \context {
    \TabStaff
    \remove Time_signature_engraver

    }
}

<<
  \new TabStaff <<\notes>>
>>
\end{lilypond}
\end{center}


\audioclip{Inicio ascendente hasta \emph{ghammaz}.}

En general, cualquiera de los inicios de esta sección puede modificarse para desembocar sobre una nota distinta de la tónica o el \emph{ghammaz}, enfatizando esta y continuando desde allí.

\subsection{Inicios en \maqam Rast}

\newbeginning{}

\begin{center}
\begin{lilypond}
\include "arabic.ly"
notes = \relative {
do16 misb32 fa misb re misb16 fa sol la sol fa sol sol sol, sol' sol do,,8
\bar "||"
}

\layout {
  \omit Voice.StringNumber
  \context {
      \TabStaff
      \tabFullNotation
      stringTunings = \stringTuning <do, fa, la, re sol do'>
    }
  indent = #0
  \context {
    \Score
    defaultBarType = ""
    supportNonIntegerFret = ##t
  }
  \context {
    \TabStaff
    \remove Time_signature_engraver

    }
}

<<
  \new TabStaff <<\notes>>
>>
\end{lilypond}
\end{center}

\audioclip{Inicio en \maqam Rast.}

\newbeginning{}

\begin{center}
\begin{lilypond}
\include "arabic.ly"
notes = \relative {
do,8 do'16 do8 do16 do8 do16 do8 do16
do,8 do'16 do8 do16 do8 do16 do8 do16
do,8 do'16 do8 sisb16 sisb8 la16 sisb8 do16
do4
\bar "||"
}

\layout {
  \omit Voice.StringNumber
  \context {
      \TabStaff
      \tabFullNotation
      stringTunings = \stringTuning <do, fa, la, re sol do'>
    }
  indent = #0
  \context {
    \Score
    defaultBarType = ""
    supportNonIntegerFret = ##t
  }
  \context {
    \TabStaff
    \remove Time_signature_engraver

    }
}

<<
  \new TabStaff <<\notes>>
>>
\end{lilypond}
\end{center}

\audioclip{Inicio en \maqam Rast.}


\newbeginning{}
El siguiente inicio arranca en la octava superior, y se utilizaría para abrir una frase intermedia, una vez que se ha llegado a esa altura en el desarrollo melódico del \taqsimb.


\begin{center}
\begin{lilypond}
\include "arabic.ly"
notes = \relative {
la16 do re do re8 re16 do do,8 do do do4
la'16 do re do re8 re16 do do,8 do do do4
\bar "||"
}

\layout {
  \omit Voice.StringNumber
  \context {
      \TabStaff
      \tabFullNotation
      stringTunings = \stringTuning <do, fa, la, re sol do'>
    }
  indent = #0
  \context {
    \Score
    defaultBarType = ""
    supportNonIntegerFret = ##t
  }
  \context {
    \TabStaff
    \remove Time_signature_engraver

    }
}

<<
  \new TabStaff <<\notes>>
>>
\end{lilypond}
\end{center}

\audioclip{Inicio en \maqam Rast.}

\newbeginning{}

Otro ejemplo similar al anterior.


\begin{center}
\begin{lilypond}
\include "arabic.ly"
notes = \relative {
r16 sol16 sol sol do do re re misb misb re re do do re32 misb re do sisb16 do do, do' do do do
\bar "||"
}

\layout {
  \omit Voice.StringNumber
  \context {
      \TabStaff
      \tabFullNotation
      stringTunings = \stringTuning <do, fa, la, re sol do'>
    }
  indent = #0
  \context {
    \Score
    defaultBarType = ""
    supportNonIntegerFret = ##t
  }
  \context {
    \TabStaff
    \remove Time_signature_engraver

    }
}

<<
  \new TabStaff <<\notes>>
>>
\end{lilypond}
\end{center}

\audioclip{Inicio en \maqam Rast.}

\newbeginning{}


\begin{center}
\begin{lilypond}
\include "arabic.ly"
notes = \relative {
sol,16 sol32 sol sol16 la sisb do re do sisb re do8 do do do
\bar "||"
}

\layout {
  \omit Voice.StringNumber
  \context {
      \TabStaff
      \tabFullNotation
      stringTunings = \stringTuning <do, fa, la, re sol do'>
    }
  indent = #0
  \context {
    \Score
    defaultBarType = ""
    supportNonIntegerFret = ##t
  }
  \context {
    \TabStaff
    \remove Time_signature_engraver

    }
}

<<
  \new TabStaff <<\notes>>
>>
\end{lilypond}
\end{center}

\audioclip{Inicio en \maqam Rast.}

\subsection{Inicios en \maqam Bayati}

\newbeginning{}

\begin{center}
\begin{lilypond}
\include "arabic.ly"
notes = \relative {
re16 sol la sib do re \tuplet 3/2 {do32 (re do)} sib16 do sib la sib \tuplet 3/2 {la32 sib la} sol16 la sol fa re
\bar "||"
}

\layout {
  \omit Voice.StringNumber
  \context {
      \TabStaff
      \tabFullNotation
      stringTunings = \stringTuning <do, fa, la, re sol do'>
    }
  indent = #0
  \context {
    \Score
    defaultBarType = ""
    supportNonIntegerFret = ##t
  }
  \context {
    \TabStaff
    \remove Time_signature_engraver

    }
}

<<
  \new TabStaff <<\notes>>
>>
\end{lilypond}
\end{center}

\audioclip{Inicio en \maqam Bayati.}

\newbeginning{}

\begin{center}
\begin{lilypond}
\include "arabic.ly"
notes = \relative {
re16 fa16 misb re do sisb do re re re re re
\bar "||"
}

\layout {
  \omit Voice.StringNumber
  \context {
      \TabStaff
      \tabFullNotation
      stringTunings = \stringTuning <do, fa, la, re sol do'>
    }
  indent = #0
  \context {
    \Score
    defaultBarType = ""
    supportNonIntegerFret = ##t
  }
  \context {
    \TabStaff
    \remove Time_signature_engraver

    }
}

<<
  \new TabStaff <<\notes>>
>>
\end{lilypond}
\end{center}

\audioclip{Inicio en \maqam Bayati.}

\newbeginning{}

\begin{center}
\begin{lilypond}
\include "arabic.ly"
notes = \relative {
re,8 re' re re, re' re re,4 do'8 re misb fa8. misb16 re8 misb re do re re re, re' re
\bar "||"
}

\layout {
  \omit Voice.StringNumber
  \context {
      \TabStaff
      \tabFullNotation
      stringTunings = \stringTuning <do, fa, la, re sol do'>
    }
  indent = #0
  \context {
    \Score
    defaultBarType = ""
    supportNonIntegerFret = ##t
  }
  \context {
    \TabStaff
    \remove Time_signature_engraver

    }
}

<<
  \new TabStaff <<\notes>>
>>
\end{lilypond}
\end{center}

\audioclip{Inicio en \maqam Bayati.}

\newbeginning{}

\begin{center}
\begin{lilypond}
\include "arabic.ly"
notes = \relative {
re16 fa8 sol la16 sib32 la sol16 fa sol sol misb fa sol sol sol
\bar "||"
}

\layout {
  \omit Voice.StringNumber
  \context {
      \TabStaff
      \tabFullNotation
      stringTunings = \stringTuning <do, fa, la, re sol do'>
    }
  indent = #0
  \context {
    \Score
    defaultBarType = ""
    supportNonIntegerFret = ##t
  }
  \context {
    \TabStaff
    \remove Time_signature_engraver

    }
}

<<
  \new TabStaff <<\notes>>
>>
\end{lilypond}
\end{center}

\audioclip{Inicio en \maqam Bayati.}

\newbeginning{}

\begin{center}
\begin{lilypond}
\include "arabic.ly"
notes = \relative {
re16 re re re re4\fermata
re16 misb fa8 fa32 misb8 misb32 fa8 fa32 misb8 misb32 fa4:32 misb16 re re4
\bar "||"
}

\layout {
  \omit Voice.StringNumber
  \context {
      \TabStaff
      \tabFullNotation
      stringTunings = \stringTuning <do, fa, la, re sol do'>
    }
  indent = #0
  \context {
    \Score
    defaultBarType = ""
    supportNonIntegerFret = ##t
  }
  \context {
    \TabStaff
    \remove Time_signature_engraver

    }
}

<<
  \new TabStaff <<\notes>>
>>
\end{lilypond}
\end{center}

\audioclip{Inicio en \maqam Bayati.}

\newbeginning{}

\begin{center}
\begin{lilypond}
\include "arabic.ly"
notes = \relative {
do16 re re re re re4\fermata
misb32 (re) do16 sisb sisb la la sol la sisb do re re re re4
\bar "||"
}

\layout {
  \omit Voice.StringNumber
  \context {
      \TabStaff
      \tabFullNotation
      stringTunings = \stringTuning <do, fa, la, re sol do'>
    }
  indent = #0
  \context {
    \Score
    defaultBarType = ""
    supportNonIntegerFret = ##t
  }
  \context {
    \TabStaff
    \remove Time_signature_engraver

    }
}

<<
  \new TabStaff <<\notes>>
>>
\end{lilypond}
\end{center}

\audioclip{Inicio en \maqam Bayati.}

\newbeginning{}

\begin{center}
\begin{lilypond}
\include "arabic.ly"
notes = \relative {
re16 re re re re re re fa32 (misb re16) re re re re re re
fa32 (misb re16) re re fa32 (misb re16) re re re re
sisb8:32 do8:32 re8:32 re16 re re4
\bar "||"
}

\layout {
  \omit Voice.StringNumber
  \context {
      \TabStaff
      \tabFullNotation
      stringTunings = \stringTuning <do, fa, la, re sol do'>
    }
  indent = #0
  \context {
    \Score
    defaultBarType = ""
    supportNonIntegerFret = ##t
  }
  \context {
    \TabStaff
    \remove Time_signature_engraver

    }
}

<<
  \new TabStaff <<\notes>>
>>
\end{lilypond}
\end{center}

\audioclip{Inicio en \maqam Bayati.}


\subsection{Inicios en \maqam Nahawand}

\newbeginning{}

\begin{center}
\begin{lilypond}
\include "arabic.ly"
notes = \relative {
do16 mib fa sol sol sol fa fa fa fa fa4
\bar "||"
}

\layout {
  \omit Voice.StringNumber
  \context {
      \TabStaff
      \tabFullNotation
      stringTunings = \stringTuning <do, fa, la, re sol do'>
    }
  indent = #0
  \context {
    \Score
    defaultBarType = ""
    supportNonIntegerFret = ##t
  }
  \context {
    \TabStaff
    \remove Time_signature_engraver

    }
}

<<
  \new TabStaff <<\notes>>
>>
\end{lilypond}
\end{center}

\audioclip{Inicio en \maqam Nahawand.}

\newbeginning{}

\begin{center}
\begin{lilypond}
\include "arabic.ly"
notes = \relative {
do8. fa16 mib8 re16 mib32 re
do8. fa16 mib8 re16 mib32 re
do8. mib16 re8 re16 do32 si
do4
\bar "||"
}

\layout {
  \omit Voice.StringNumber
  \context {
      \TabStaff
      \tabFullNotation
      stringTunings = \stringTuning <do, fa, la, re sol do'>
    }
  indent = #0
  \context {
    \Score
    defaultBarType = ""
    supportNonIntegerFret = ##t
  }
  \context {
    \TabStaff
    \remove Time_signature_engraver

    }
}

<<
  \new TabStaff <<\notes>>
>>
\end{lilypond}
\end{center}

\audioclip{Inicio en \maqam Nahawand.}


\newbeginning{}

\begin{center}
\begin{lilypond}
\include "arabic.ly"
notes = \relative {
do,16 do' do do do si (do4)
do16 re mib fa, mib'16 (fa4)
fa16 mib mib re fa mib mib re re
do do si (do4)
\bar "||"
}

\layout {
  \omit Voice.StringNumber
  \context {
      \TabStaff
      \tabFullNotation
      stringTunings = \stringTuning <do, fa, la, re sol do'>
    }
  indent = #0
  \context {
    \Score
    defaultBarType = ""
    supportNonIntegerFret = ##t
  }
  \context {
    \TabStaff
    \remove Time_signature_engraver

    }
}

<<
  \new TabStaff <<\notes>>
>>
\end{lilypond}
\end{center}

\audioclip{Inicio en \maqam Nahawand.}

\newbeginning{}

\begin{center}
\begin{lilypond}
\include "arabic.ly"
notes = \relative {
do'16 re\glissando mib8^\prall re^\prall do sib^\prall do16 do,, do''4
\bar "||"
}

\layout {
  \omit Voice.StringNumber
  \context {
      \TabStaff
      \tabFullNotation
      stringTunings = \stringTuning <do, fa, la, re sol do'>
    }
  indent = #0
  \context {
    \Score
    defaultBarType = ""
    supportNonIntegerFret = ##t
  }
  \context {
    \TabStaff
    \remove Time_signature_engraver

    }
}

<<
  \new TabStaff <<\notes>>
>>
\end{lilypond}
\end{center}

\audioclip{Inicio en \maqam Nahawand.}


\subsection{Inicios en \maqam Kurd}

\newbeginning{}

\begin{center}
\begin{lilypond}
\include "arabic.ly"
notes = \relative {
do16 re mib fa fa fa fa fa mib re do re8 re re
\bar "||"
}

\layout {
  \omit Voice.StringNumber
  \context {
      \TabStaff
      \tabFullNotation
      stringTunings = \stringTuning <do, fa, la, re sol do'>
    }
  indent = #0
  \context {
    \Score
    defaultBarType = ""
    supportNonIntegerFret = ##t
  }
  \context {
    \TabStaff
    \remove Time_signature_engraver

    }
}

<<
  \new TabStaff <<\notes>>
>>
\end{lilypond}
\end{center}

\audioclip{Inicio en \maqam Kurd.}

\newbeginning{}

\begin{center}
\begin{lilypond}
\include "arabic.ly"
notes = \relative {
re8 do re mib fa4 mib8 re re4
\bar "||"
}

\layout {
  \omit Voice.StringNumber
  \context {
      \TabStaff
      \tabFullNotation
      stringTunings = \stringTuning <do, fa, la, re sol do'>
    }
  indent = #0
  \context {
    \Score
    defaultBarType = ""
    supportNonIntegerFret = ##t
  }
  \context {
    \TabStaff
    \remove Time_signature_engraver

    }
}

<<
  \new TabStaff <<\notes>>
>>
\end{lilypond}
\end{center}

\audioclip{Inicio en \maqam Kurd.}


\subsection{Inicios en \maqam Hiyaz}

\newbeginning{}

\begin{center}
\begin{lilypond}
\include "arabic.ly"
notes = \relative {
\time 3/4
r16 do16 re re re re do, (re) re'16 re re re re4
do16 re mib4\glissando fad4\glissando mib8 re16 re, re' re re
\bar "||"
}

\layout {
  \omit Voice.StringNumber
  \context {
      \TabStaff
      \tabFullNotation
      stringTunings = \stringTuning <do, fa, la, re sol do'>
    }
  indent = #0
  \context {
    \Score
    defaultBarType = ""
    supportNonIntegerFret = ##t
  }
  \context {
    \TabStaff
    \remove Time_signature_engraver

    }
}

<<
  \new TabStaff <<\notes>>
>>
\end{lilypond}
\end{center}

\audioclip{Inicio en \maqam Hiyaz.}


\section{Cierres}

Esta sección contiene fragmentos empleados para concluir una frase o un \taqsimb, también conocidos como \emph{qafla}.

Comencemos con una de las opciones más habituales: un descenso a tónica empleando un motivo melódico. Por ejemplo:


\begin{center}
\begin{lilypond}
\include "arabic.ly"
notes = \relative {
\key do \rast
la8 sol fa misb sol fa misb re fa misb \tuplet 3/2 {re16 (misb re)} do4
\bar "||"
}

\layout {
  \omit Voice.StringNumber
  \context {
      \TabStaff
      \tabFullNotation
      stringTunings = \stringTuning <do, fa, la, re sol do'>
    }
  indent = #0
  \context {
    \Score
    defaultBarType = ""
    supportNonIntegerFret = ##t
  }
  \context {
    \TabStaff
    \remove Time_signature_engraver

    }
}

<<
  \new TabStaff <<\notes>>
>>
\end{lilypond}
\end{center}


\audioclip{Cierre con descenso y motivo melódico.}

Es habitual que el descenso a la tónica se produzca desde las cercanías del \emph{ghammaz} o incluso desde la octava, pero si la frase se encuentra cerca de la tónica, puede saltarse hasta una nota superior y comenzar desde allí. Este es uno de los casos en los que aparecen intervalos amplios dentro de la melodía. En el siguiente ejemplo, en el \maqam Nahawand, se salta desde la tónica al \emph{ghammaz} y desde allí se desciende mediante un patrón melodico, al que sigue un cierre con trémolo.


\begin{center}
\begin{lilypond}
\include "arabic.ly"
notes = \relative {
\key do \minor
do8 do do mib re do
\tuplet 3/2 {sol' fa mib} \tuplet 3/2 {fa mib re} \tuplet 3/2  {mib re do} do4 \tuplet 3/2 {do8 re mib}
fa4:32 re16 do
\bar "||"
}

\layout {
  \omit Voice.StringNumber
  \context {
      \TabStaff
      \tabFullNotation
      stringTunings = \stringTuning <do, fa, la, re sol do'>
    }
  indent = #0
  \context {
    \Score
    defaultBarType = ""
    supportNonIntegerFret = ##t
  }
  \context {
    \TabStaff
    \remove Time_signature_engraver

    }
}

<<
  \new TabStaff <<\notes>>
>>
\end{lilypond}
\end{center}


\audioclip{Salto hasta \emph{ghammaz} y cierre con descenso.}

El uso del trémolo para enfatizar un cierre es casi obligado en algunos estilos de ud como el árabe. El cierre más simple con uso de trémolo es un simple descenso a tónica siguiendo las notas del \maqamb.


\begin{center}
\begin{lilypond}
\include "arabic.ly"
notes = \relative {
\key do \rast
fa4:32 misb4:32 re4:32 do2
\bar "||"
}

\layout {
  \omit Voice.StringNumber
  \context {
      \TabStaff
      \tabFullNotation
      stringTunings = \stringTuning <do, fa, la, re sol do'>
    }
  indent = #0
  \context {
    \Score
    defaultBarType = ""
    supportNonIntegerFret = ##t
  }
  \context {
    \TabStaff
    \remove Time_signature_engraver

    }
}

<<
  \new TabStaff <<\notes>>
>>
\end{lilypond}
\end{center}


\audioclip{Cierre simple con trémolo.}

El siguiente cierre combina el trémolo con un motivo melódico:


\begin{center}
\begin{lilypond}
\include "arabic.ly"
notes = \relative {
\key do \rast
la8:64 [sol8:64 fa8:64] sol8:64 [fa8:64 misb8:64] fa8:64 [misb8:64 re8:64] do4
\bar "||"
}

\layout {
  \omit Voice.StringNumber
  \context {
      \TabStaff
      \tabFullNotation
      stringTunings = \stringTuning <do, fa, la, re sol do'>
    }
  indent = #0
  \context {
    \Score
    defaultBarType = ""
    supportNonIntegerFret = ##t
  }
  \context {
    \TabStaff
    \remove Time_signature_engraver

    }
}

<<
  \new TabStaff <<\notes>>
>>
\end{lilypond}
\end{center}


\audioclip{Cierre con trémolo y motivo melódico.}

El uso de la tónica una octava por debajo es un recurso frecuente para dar solidez a un remate.


\begin{center}
\begin{lilypond}
\include "arabic.ly"
notes = \relative {
\key do \rast
sol8 (la16 sol) fa4 fa8 (sol16\3 fa) misb4 misb8 (fa16 misb) re4 re8 (misb16 re) do4
do,2
\bar "||"
}

\layout {
  \omit Voice.StringNumber
  \context {
      \TabStaff
      \tabFullNotation
      stringTunings = \stringTuning <do, fa, la, re sol do'>
    }
  indent = #0
  \context {
    \Score
    defaultBarType = ""
    supportNonIntegerFret = ##t
  }
  \context {
    \TabStaff
    \remove Time_signature_engraver

    }
}

<<
  \new TabStaff <<\notes>>
>>
\end{lilypond}
\end{center}


\audioclip{Cierre con tónica en octava inferior.}

El siguiente ejemplo en \maqam Bayati usa octavas y trémolo conjuntamente:


\begin{center}
\begin{lilypond}
\include "arabic.ly"
notes = \relative {
\key re \bayati
[re16 misb fa] [sol fa sol] fa,16 fa'4:32 sol,16 sol'4:32 misb8.:64 re32 do re4
\bar "||"
}

\layout {
  \omit Voice.StringNumber
  \context {
      \TabStaff
      \tabFullNotation
      stringTunings = \stringTuning <do, fa, la, re sol do'>
    }
  indent = #0
  \context {
    \Score
    defaultBarType = ""
    supportNonIntegerFret = ##t
  }
  \context {
    \TabStaff
    \remove Time_signature_engraver

    }
}

<<
  \new TabStaff <<\notes>>
>>
\end{lilypond}
\end{center}


\audioclip{Cierre con octavas y trémolo.}

Aun si el patrón es descendente, muchos cierres usan la nota inmediatamente inferior a la tónica como nota sensible desde la que resolver y concluir la \emph{qafla}, en dirección ascendente. El siguiente es un ejemplo de ello:


\begin{center}
\begin{lilypond}
\include "arabic.ly"
notes = \relative {
\key do \rast
do8 re misb fa fa fa fa misb16 fa misb re misb8 re do sisb do4
\bar "||"
}

\layout {
  \omit Voice.StringNumber
  \context {
      \TabStaff
      \tabFullNotation
      stringTunings = \stringTuning <do, fa, la, re sol do'>
    }
  indent = #0
  \context {
    \Score
    defaultBarType = ""
    supportNonIntegerFret = ##t
  }
  \context {
    \TabStaff
    \remove Time_signature_engraver

    }
}

<<
  \new TabStaff <<\notes>>
>>
\end{lilypond}
\end{center}


\audioclip{Cierre con uso de nota sensible.}

Aunque menos frecuente, también es posible apoyarse en el \emph{ghammaz} como nota previa antes de la tónica, que luego resuelve sobre esta. El siguiente ejemplo combina esto con el uso de trémolo y octavas.


\begin{center}
\begin{lilypond}
\include "arabic.ly"
notes = \relative {
\key do \rast
fa,8 fa'4.:32 sol,8 sol'4.:32 do,,8 do'2
\bar "||"
}

\layout {
  \omit Voice.StringNumber
  \context {
      \TabStaff
      \tabFullNotation
      stringTunings = \stringTuning <do, fa, la, re sol do'>
    }
  indent = #0
  \context {
    \Score
    defaultBarType = ""
    supportNonIntegerFret = ##t
  }
  \context {
    \TabStaff
    \remove Time_signature_engraver

    }
}

<<
  \new TabStaff <<\notes>>
>>
\end{lilypond}
\end{center}


\audioclip{Cierre con uso de \emph{ghammaz} en octava inferior.}

La nota inmediatamente por encima de la tónica, especialmente si se encuentra a un semitono de aquella (como sucede en el \maqam Kurd, empleado en el siguiente ejemplo), puede también utilizarse con carácter de nota sensible, deteniéndose en ella y enfatizándola antes de resolver.


\begin{center}
\begin{lilypond}
\include "arabic.ly"
notes = \relative {
\key do \rast
\time 16/4
mib16 fa sol la sol fa re mib4:64\fermata (re32) do re16 re, re'2
\bar "||"
}

\layout {
  \omit Voice.StringNumber
  \context {
      \TabStaff
      \tabFullNotation
      stringTunings = \stringTuning <do, fa, la, re sol do'>
    }
  indent = #0
  \context {
    \Score
    defaultBarType = ""
    supportNonIntegerFret = ##t
  }
  \context {
    \TabStaff
    \remove Time_signature_engraver

    }
}

<<
  \new TabStaff <<\notes>>
>>
\end{lilypond}
\end{center}


\audioclip{Cierre con nota inmediatamente superior a la tónica.}

Ampliando esta idea, otra alternativa habitual es llegar a la tónica mediante un motivo melódico ascendente, desde la octava inferior. El siguiente ejemplo combina un patrón melódico descendente que baja a la octava inferior, con otro ascendente que resuelve desde allí sobre la tónica.


\begin{center}
\begin{lilypond}
\include "arabic.ly"
notes = \relative {
\key do \rast
fa8:32 misb8:32 re8:32 do8:32 sisb8:32 la8:32 sol4 sol16 la sib do2
\bar "||"
}

\layout {
  \omit Voice.StringNumber
  \context {
      \TabStaff
      \tabFullNotation
      stringTunings = \stringTuning <do, fa, la, re sol do'>
    }
  indent = #0
  \context {
    \Score
    defaultBarType = ""
    supportNonIntegerFret = ##t
  }
  \context {
    \TabStaff
    \remove Time_signature_engraver

    }
}

<<
  \new TabStaff <<\notes>>
>>
\end{lilypond}
\end{center}


\audioclip{Cierre con patrón ascendente.}

En ocasiones, la \emph{qafla} tiene dos partes, o más bien se puede decir que se unen dos remates consecutivos, ambos acabando sobre la tónica. Por ejemplo:


\begin{center}
\begin{lilypond}
\include "arabic.ly"
notes = \relative {
\key do \rast
sol8 (la16 sol) fa4 fa8 (sol16\3 fa) misb4 misb8 (fa16 misb) re4 re8 (misb16 re) do4
do8 re misb fa4 fa4:32 misb4 misb4:32 re4 re4:32 do
\bar "||"
}

\layout {
  \omit Voice.StringNumber
  \context {
      \TabStaff
      \tabFullNotation
      stringTunings = \stringTuning <do, fa, la, re sol do'>
    }
  indent = #0
  \context {
    \Score
    defaultBarType = ""
    supportNonIntegerFret = ##t
  }
  \context {
    \TabStaff
    \remove Time_signature_engraver

    }
}

<<
  \new TabStaff <<\notes>>
>>
\end{lilypond}
\end{center}


\audioclip{Dos cierres consecutivos.}

Otro ejemplo similar, en este caso en \maqam Nahawand:


\begin{center}
\begin{lilypond}
\include "arabic.ly"
notes = \relative {
\key do \minor
sol16-> sol sol sol
fa-> fa fa fa
mib-> mib mib mib
re-> re re re
do8
do16 re re mib mib fa fa sol sol
la sol fa mib
sol fa mib re
do
\bar "||"
}

\layout {
  \omit Voice.StringNumber
  \context {
      \TabStaff
      \tabFullNotation
      stringTunings = \stringTuning <do, fa, la, re sol do'>
    }
  indent = #0
  \context {
    \Score
    defaultBarType = ""
    supportNonIntegerFret = ##t
  }
  \context {
    \TabStaff
    \remove Time_signature_engraver

    }
}

<<
  \new TabStaff <<\notes>>
>>
\end{lilypond}
\end{center}



La siguiente frase se emplea con frecuencia como segundo cierre, aunque su caracter es más de ornamentación, y es el cierre previo el que deja la melodía concluida sobre la tónica. Se suele tocar con menos volumen que el de la frase anterior, y es habitual no en el cierre del \taqsimb, sino en el de frases intermedias.


\begin{center}
\begin{lilypond}
\include "arabic.ly"
notes = \relative {
\key do \rast
r8 \tuplet 3/2 {re16 (misb fa)} sol4 \tuplet 3/2 {fa16 (misb re)} do4
\bar "||"
}

\layout {
  \omit Voice.StringNumber
  \context {
      \TabStaff
      \tabFullNotation
      stringTunings = \stringTuning <do, fa, la, re sol do'>
    }
  indent = #0
  \context {
    \Score
    defaultBarType = ""
    supportNonIntegerFret = ##t
  }
  \context {
    \TabStaff
    \remove Time_signature_engraver

    }
}

<<
  \new TabStaff <<\notes>>
>>
\end{lilypond}
\end{center}


\audioclip{Frase de doble cierre.}

Otro ejemplo más de esta idea:


\begin{center}
\begin{lilypond}
\include "arabic.ly"
notes = \relative {
\key do \rast
fa8 fa misb16 fa misb re misb16 re16 do sisb do4
misb16 (re) misb16 (re) misb16 (re) do sisb do4
\bar "||"
}

\layout {
  \omit Voice.StringNumber
  \context {
      \TabStaff
      \tabFullNotation
      stringTunings = \stringTuning <do, fa, la, re sol do'>
    }
  indent = #0
  \context {
    \Score
    defaultBarType = ""
    supportNonIntegerFret = ##t
  }
  \context {
    \TabStaff
    \remove Time_signature_engraver

    }
}

<<
  \new TabStaff <<\notes>>
>>
\end{lilypond}
\end{center}


\audioclip{Frase de doble cierre.}

Respecto a la duración de la última nota del remate, lo habitual es dejar esta sonando. Sin embargo, en el cierre de frases interiores de un \taqsim pueden emplearse finales secos, tras los que normalmente comienza rápido la siguiente frase. Estos cierres pueden ocurrir sobre la tónica, pero también sobre otras notas del \maqamb. Por ejemplo, la siguiente frase en \maqam Nahawand se cierra sobre la nota re.


\begin{center}
\begin{lilypond}
\include "arabic.ly"
notes = \relative {
r8 do8 re mib fa sol fa sol fa4 mib32 (fa mib re) r4
\bar "||"
}

\layout {
  \omit Voice.StringNumber
  \context {
      \TabStaff
      \tabFullNotation
      stringTunings = \stringTuning <do, fa, la, re sol do'>
    }
  indent = #0
  \context {
    \Score
    defaultBarType = ""
    supportNonIntegerFret = ##t
  }
  \context {
    \TabStaff
    \remove Time_signature_engraver

    }
}

<<
  \new TabStaff <<\notes>>
>>
\end{lilypond}
\end{center}


\audioclip{Cierres con final seco.}

En el estilo árabe, como ya se ha dicho, el cierre que da fin al \taqsim ha de ser intenso, con un cierto caracter de apoteosis. Para ello, es necesario que exista una tensión melódica antes de llegar al remate, y esta se debe crear progresivamente en las frases previas. Una forma habitual de hacer esto es recorrer el \maqam de manera ascendente con algún patrón melódico, para despues descender y concluir sobre la tónica. Por ejemplo:


\begin{center}
\begin{lilypond}
\include "arabic.ly"
notes = \relative {
\key do \rast
do16 re misb8 re16 misb fa8
misb16 fa sol8 fa16 sol la8
sol16 la sisb8 la16 sisb do4
\bar "|"
sisb16 do la sisb sol la fa sol
misb fa re misb fa4 fa8:32 misb4 misb8:32 re4 re8:32 do4
\bar "||"
}

\layout {
  \omit Voice.StringNumber
  \context {
      \TabStaff
      \tabFullNotation
      stringTunings = \stringTuning <do, fa, la, re sol do'>
    }
  indent = #0
  \context {
    \Score
    defaultBarType = ""
    supportNonIntegerFret = ##t
  }
  \context {
    \TabStaff
    \remove Time_signature_engraver

    }
}

<<
  \new TabStaff <<\notes>>
>>
\end{lilypond}
\end{center}


\audioclip{Cierre con ascenso previo para crear tensión.}

Otra posibilidad sobre esta misma idea:


\begin{center}
\begin{lilypond}
\include "arabic.ly"
notes = \relative {
\key do \rast
do8 re misb re misb fa misb fa sol fa sol la4
sisb8 sisb la la sol sol fa fa misb misb fa4:32
misb32 (fa misb re) do4
\bar "||"
}

\layout {
  \omit Voice.StringNumber
  \context {
      \TabStaff
      \tabFullNotation
      stringTunings = \stringTuning <do, fa, la, re sol do'>
    }
  indent = #0
  \context {
    \Score
    defaultBarType = ""
    supportNonIntegerFret = ##t
  }
  \context {
    \TabStaff
    \remove Time_signature_engraver

    }
}

<<
  \new TabStaff <<\notes>>
>>
\end{lilypond}
\end{center}

\audioclip{Cierre con ascenso previo para crear tensión.}

La sección ascendente es frecuente tocarla acelerando el tempo, para así añadir más tensión al desarrollo.

\subsection{Cierres en \maqam Rast}

\newending{}

\begin{center}
\begin{lilypond}
\include "arabic.ly"
notes = \relative {
\time 3/4
fa8:32 misb8:32 re8:32
fa8:32 misb8:32 re8:32
re16 do16 do sisb do4
\bar "||"
}

\layout {
  \omit Voice.StringNumber
  \context {
      \TabStaff
      \tabFullNotation
      stringTunings = \stringTuning <do, fa, la, re sol do'>
    }
  indent = #0
  \context {
    \Score
    defaultBarType = ""
    supportNonIntegerFret = ##t
  }
  \context {
    \TabStaff
    \remove Time_signature_engraver

    }
}

<<
  \new TabStaff <<\notes>>
>>
\end{lilypond}
\end{center}

\audioclip{Cierre en \maqam Rast.}

\newending{}

\begin{center}
\begin{lilypond}
\include "arabic.ly"
notes = \relative {
do'16 sisb sisb
do sisb sisb la8.
sol32 (la sol) fa8. misb32 do8
\bar "||"
}

\layout {
  \omit Voice.StringNumber
  \context {
      \TabStaff
      \tabFullNotation
      stringTunings = \stringTuning <do, fa, la, re sol do'>
    }
  indent = #0
  \context {
    \Score
    defaultBarType = ""
    supportNonIntegerFret = ##t
  }
  \context {
    \TabStaff
    \remove Time_signature_engraver

    }
}

<<
  \new TabStaff <<\notes>>
>>
\end{lilypond}
\end{center}

\audioclip{Cierre en \maqam Rast.}

\newending{}

\begin{center}
\begin{lilypond}
\include "arabic.ly"
notes = \relative {
fa16 fa fa fa
sol16 fa misb32 (fa) re (misb) do16 do do, do'4
\bar "||"
}

\layout {
  \omit Voice.StringNumber
  \context {
      \TabStaff
      \tabFullNotation
      stringTunings = \stringTuning <do, fa, la, re sol do'>
    }
  indent = #0
  \context {
    \Score
    defaultBarType = ""
    supportNonIntegerFret = ##t
  }
  \context {
    \TabStaff
    \remove Time_signature_engraver

    }
}

<<
  \new TabStaff <<\notes>>
>>
\end{lilypond}
\end{center}

\audioclip{Cierre en \maqam Rast.}

\newending{}

\begin{center}
\begin{lilypond}
\include "arabic.ly"
notes = \relative {
\tuplet 3/2 {misb16 (fa) sol}
\tuplet 3/2 {misb16 (fa) sol}
\tuplet 3/2 {misb16 (fa) sol}
fa4:32 misb32 (re) do4
\bar "||"
}

\layout {
  \omit Voice.StringNumber
  \context {
      \TabStaff
      \tabFullNotation
      stringTunings = \stringTuning <do, fa, la, re sol do'>
    }
  indent = #0
  \context {
    \Score
    defaultBarType = ""
    supportNonIntegerFret = ##t
  }
  \context {
    \TabStaff
    \remove Time_signature_engraver

    }
}

<<
  \new TabStaff <<\notes>>
>>
\end{lilypond}
\end{center}

\audioclip{Cierre en \maqam Rast.}

\newending{}

El cierre sobre la tónica se da habitualmente sobre el do central del ud, y el do más grave (sexta cuerda al aire) se usa para puntuar aquella. Sin embargo, es posible utilizar toda la extensión del maqam en la octava inferior para cerrar sobre este último, como en el siguiente ejemplo.


\begin{center}
\begin{lilypond}
\include "arabic.ly"
notes = \relative {
do'8 sol
misb fa sol fa misb re do re do sisb la sol fa misb re do4
\bar "||"
}

\layout {
  \omit Voice.StringNumber
  \context {
      \TabStaff
      \tabFullNotation
      stringTunings = \stringTuning <do, fa, la, re sol do'>
    }
  indent = #0
  \context {
    \Score
    defaultBarType = ""
    supportNonIntegerFret = ##t
  }
  \context {
    \TabStaff
    \remove Time_signature_engraver

    }
}

<<
  \new TabStaff <<\notes>>
>>
\end{lilypond}
\end{center}

\audioclip{Cierre en \maqam Rast.}

\newending{}

\begin{center}
\begin{lilypond}
\include "arabic.ly"
notes = \relative {
sol16 la sisb do4
sol,16  la sisb do4
\bar "||"
}

\layout {
  \omit Voice.StringNumber
  \context {
      \TabStaff
      \tabFullNotation
      stringTunings = \stringTuning <do, fa, la, re sol do'>
    }
  indent = #0
  \context {
    \Score
    defaultBarType = ""
    supportNonIntegerFret = ##t
  }
  \context {
    \TabStaff
    \remove Time_signature_engraver

    }
}

<<
  \new TabStaff <<\notes>>
>>
\end{lilypond}
\end{center}

\audioclip{Cierre en \maqam Rast.}

\newending{}

\begin{center}
\begin{lilypond}
\include "arabic.ly"
notes = \relative {
fa16 fa misb fa fa misb fa8:64 misb16 fa sol la sol fa misb re do
\bar "||"
}

\layout {
  \omit Voice.StringNumber
  \context {
      \TabStaff
      \tabFullNotation
      stringTunings = \stringTuning <do, fa, la, re sol do'>
    }
  indent = #0
  \context {
    \Score
    defaultBarType = ""
    supportNonIntegerFret = ##t
  }
  \context {
    \TabStaff
    \remove Time_signature_engraver

    }
}

<<
  \new TabStaff <<\notes>>
>>
\end{lilypond}
\end{center}

\audioclip{Cierre en \maqam Rast.}

\newending{}

\begin{center}
\begin{lilypond}
\include "arabic.ly"
notes = \relative {
sisb,8 re8:64 do8 misb8:64 re8 fa8:64 misb8 sol8:64 fa4
sol8 sol8:64 fa8 fa8:64 misb8 misb8:64 re8 re8:64 do4
\bar "||"
}

\layout {
  \omit Voice.StringNumber
  \context {
      \TabStaff
      \tabFullNotation
      stringTunings = \stringTuning <do, fa, la, re sol do'>
    }
  indent = #0
  \context {
    \Score
    defaultBarType = ""
    supportNonIntegerFret = ##t
  }
  \context {
    \TabStaff
    \remove Time_signature_engraver

    }
}

<<
  \new TabStaff <<\notes>>
>>
\end{lilypond}
\end{center}

\audioclip{Cierre en \maqam Rast.}

\newending{}

\begin{center}
\begin{lilypond}
\include "arabic.ly"
notes = \relative {
sol16 la32 sol sol16 do8:32 la16 sol fa misb re do4
\bar "||"
}

\layout {
  \omit Voice.StringNumber
  \context {
      \TabStaff
      \tabFullNotation
      stringTunings = \stringTuning <do, fa, la, re sol do'>
    }
  indent = #0
  \context {
    \Score
    defaultBarType = ""
    supportNonIntegerFret = ##t
  }
  \context {
    \TabStaff
    \remove Time_signature_engraver

    }
}

<<
  \new TabStaff <<\notes>>
>>
\end{lilypond}
\end{center}

\audioclip{Cierre en \maqam Rast.}

\newending{}

\begin{center}
\begin{lilypond}
\include "arabic.ly"
notes = \relative {
fa8. sol16 fa8 sol
misb8. fa16 misb8 fa
re8 re16 misb re8 misb do2
fa16 misb re do4
\bar "||"
}

\layout {
  \omit Voice.StringNumber
  \context {
      \TabStaff
      \tabFullNotation
      stringTunings = \stringTuning <do, fa, la, re sol do'>
    }
  indent = #0
  \context {
    \Score
    defaultBarType = ""
    supportNonIntegerFret = ##t
  }
  \context {
    \TabStaff
    \remove Time_signature_engraver

    }
}

<<
  \new TabStaff <<\notes>>
>>
\end{lilypond}
\end{center}

\audioclip{Cierre en \maqam Rast.}

\subsection{Cierres en \maqam Bayati}

\newending{}

\begin{center}
\begin{lilypond}
\include "arabic.ly"
notes = \relative {
r16 sib16 (la8) la16  (sol8) sol16 fa8 sol4 sol16 fa (misb) misb4:32 (re16) re, re'4
\bar "||"
}

\layout {
  \omit Voice.StringNumber
  \context {
      \TabStaff
      \tabFullNotation
      stringTunings = \stringTuning <do, fa, la, re sol do'>
    }
  indent = #0
  \context {
    \Score
    defaultBarType = ""
    supportNonIntegerFret = ##t
  }
  \context {
    \TabStaff
    \remove Time_signature_engraver

    }
}

<<
  \new TabStaff <<\notes>>
>>
\end{lilypond}
\end{center}


\audioclip{Cierre en \maqam Bayati.}

\newending{}

\begin{center}
\begin{lilypond}
\include "arabic.ly"
notes = \relative {
sol16 sol fa fa misb misb fa4:32 misb32 (fa misb re) do16 re8 re4
\bar "||"
}

\layout {
  \omit Voice.StringNumber
  \context {
      \TabStaff
      \tabFullNotation
      stringTunings = \stringTuning <do, fa, la, re sol do'>
    }
  indent = #0
  \context {
    \Score
    defaultBarType = ""
    supportNonIntegerFret = ##t
  }
  \context {
    \TabStaff
    \remove Time_signature_engraver

    }
}

<<
  \new TabStaff <<\notes>>
>>
\end{lilypond}
\end{center}


\audioclip{Cierre en \maqam Bayati.}

\newending{}

\begin{center}
\begin{lilypond}
\include "arabic.ly"
notes = \relative {
re8 misb fa sol fa misb fa4:32 do,16 (re8.)
\bar "||"
}

\layout {
  \omit Voice.StringNumber
  \context {
      \TabStaff
      \tabFullNotation
      stringTunings = \stringTuning <do, fa, la, re sol do'>
    }
  indent = #0
  \context {
    \Score
    defaultBarType = ""
    supportNonIntegerFret = ##t
  }
  \context {
    \TabStaff
    \remove Time_signature_engraver

    }
}

<<
  \new TabStaff <<\notes>>
>>
\end{lilypond}
\end{center}


\audioclip{Cierre en \maqam Bayati.}

\newending{}

\begin{center}
\begin{lilypond}
\include "arabic.ly"
notes = \relative {
\partial 8.
misb16 fa sol la8 sol16 fa8
re16 misb fa sol8 fa16 misb8
do16 re misb fa4 misb16 re4 do16 sisb do re re4
\bar "||"
}

\layout {
  \omit Voice.StringNumber
  \context {
      \TabStaff
      \tabFullNotation
      stringTunings = \stringTuning <do, fa, la, re sol do'>
    }
  indent = #0
  \context {
    \Score
    defaultBarType = ""
    supportNonIntegerFret = ##t
  }
  \context {
    \TabStaff
    \remove Time_signature_engraver

    }
}

<<
  \new TabStaff <<\notes>>
>>
\end{lilypond}
\end{center}


\audioclip{Cierre en \maqam Bayati.}

\newending{}

\begin{center}
\begin{lilypond}
\include "arabic.ly"
notes = \relative {
misb4:32 fa4:32 fa16 misb fa sol re4 misb32 (re) do16 re re4
\bar "||"
}

\layout {
  \omit Voice.StringNumber
  \context {
      \TabStaff
      \tabFullNotation
      stringTunings = \stringTuning <do, fa, la, re sol do'>
    }
  indent = #0
  \context {
    \Score
    defaultBarType = ""
    supportNonIntegerFret = ##t
  }
  \context {
    \TabStaff
    \remove Time_signature_engraver

    }
}

<<
  \new TabStaff <<\notes>>
>>
\end{lilypond}
\end{center}


\audioclip{Cierre en \maqam Bayati.}

\audioclip{Cierre en \maqam Bayati.}

\newending{}

\begin{center}
\begin{lilypond}
\include "arabic.ly"
notes = \relative {
sol16 la fa sol sol sol fa4:32 (re16) re, re'4
\bar "||"
}

\layout {
  \omit Voice.StringNumber
  \context {
      \TabStaff
      \tabFullNotation
      stringTunings = \stringTuning <do, fa, la, re sol do'>
    }
  indent = #0
  \context {
    \Score
    defaultBarType = ""
    supportNonIntegerFret = ##t
  }
  \context {
    \TabStaff
    \remove Time_signature_engraver

    }
}

<<
  \new TabStaff <<\notes>>
>>
\end{lilypond}
\end{center}


\audioclip{Cierre en \maqam Bayati.}

\subsection{Cierres en \maqam Nahawand}

\newending{}
%simon shaheen

\begin{center}
\begin{lilypond}
\include "arabic.ly"
notes = \relative {
sol8 fa sol mib16. (re32) do8 do16 re mib fa sol32 (fa) mib16 fa32 (mib) re16 mib32 (re) do8
\tuplet 3/2 {re16 (mib fa)} mib (fa) mib8 mib32 (re) do8
\bar "||"
}

\layout {
  \omit Voice.StringNumber
  \context {
      \TabStaff
      \tabFullNotation
      stringTunings = \stringTuning <do, fa, la, re sol do'>
    }
  indent = #0
  \context {
    \Score
    defaultBarType = ""
    supportNonIntegerFret = ##t
  }
  \context {
    \TabStaff
    \remove Time_signature_engraver

    }
}

<<
  \new TabStaff <<\notes>>
>>
\end{lilypond}
\end{center}

\audioclip{Cierre en \maqam Nahawand.}

\newending{}
%Omar Abbad

\begin{center}
\begin{lilypond}
\include "arabic.ly"
notes = \relative {
do'4 do16 sib8. sib16 lab8. lab16 sol8. sol16 fa8. fa16 mib8. mib16 re8. re16 do8.
re16 mib fa8 mib sol2:32 sol16 fa mib re do
\bar "||"
}

\layout {
  \omit Voice.StringNumber
  \context {
      \TabStaff
      \tabFullNotation
      stringTunings = \stringTuning <do, fa, la, re sol do'>
    }
  indent = #0
  \context {
    \Score
    defaultBarType = ""
    supportNonIntegerFret = ##t
  }
  \context {
    \TabStaff
    \remove Time_signature_engraver

    }
}

<<
  \new TabStaff <<\notes>>
>>
\end{lilypond}
\end{center}

\audioclip{Cierre en \maqam Nahawand.}

\newending{}

\begin{center}
\begin{lilypond}
\include "arabic.ly"
notes = \relative {
\tuplet 3/2 {do8 sol do} \tuplet 3/2 {mib do mib} \tuplet 3/2 {sol mib sol} do4\2\prall
\bar "||"
}

\layout {
  \omit Voice.StringNumber
  \context {
      \TabStaff
      \tabFullNotation
      stringTunings = \stringTuning <do, fa, la, re sol do'>
    }
  indent = #0
  \context {
    \Score
    defaultBarType = ""
    supportNonIntegerFret = ##t
  }
  \context {
    \TabStaff
    \remove Time_signature_engraver

    }
}

<<
  \new TabStaff <<\notes>>
>>
\end{lilypond}
\end{center}

\audioclip{Cierre en \maqam Nahawand.}

\newending{}

\begin{center}
\begin{lilypond}
\include "arabic.ly"
notes = \relative {
sol8 fa sol lab sol16 sol fa fa mib mib re mib32 re do8 re mib32 fa mib16 mib4:64  (re32) do4
\bar "||"
}

\layout {
  \omit Voice.StringNumber
  \context {
      \TabStaff
      \tabFullNotation
      stringTunings = \stringTuning <do, fa, la, re sol do'>
    }
  indent = #0
  \context {
    \Score
    defaultBarType = ""
    supportNonIntegerFret = ##t
  }
  \context {
    \TabStaff
    \remove Time_signature_engraver

    }
}

<<
  \new TabStaff <<\notes>>
>>
\end{lilypond}
\end{center}

\audioclip{Cierre en \maqam Nahawand.}

\newending{}

\begin{center}
\begin{lilypond}
\include "arabic.ly"
notes = \relative {
sol8 fa32 mib16 fa8 mib32 re16 mib8 re32 do16 re do8
lab'16 sol fa sol mib re do4
\bar "||"
}

\layout {
  \omit Voice.StringNumber
  \context {
      \TabStaff
      \tabFullNotation
      stringTunings = \stringTuning <do, fa, la, re sol do'>
    }
  indent = #0
  \context {
    \Score
    defaultBarType = ""
    supportNonIntegerFret = ##t
  }
  \context {
    \TabStaff
    \remove Time_signature_engraver

    }
}

<<
  \new TabStaff <<\notes>>
>>
\end{lilypond}
\end{center}

\audioclip{Cierre en \maqam Nahawand.}

\newending{}

\begin{center}
\begin{lilypond}
\include "arabic.ly"
notes = \relative {
sol8 sol8:32 lab16 (sol) fa4 mib16 (fa) sol fa (mib re) do (re mib re) do (si) do4
\bar "||"
}

\layout {
  \omit Voice.StringNumber
  \context {
      \TabStaff
      \tabFullNotation
      stringTunings = \stringTuning <do, fa, la, re sol do'>
    }
  indent = #0
  \context {
    \Score
    defaultBarType = ""
    supportNonIntegerFret = ##t
  }
  \context {
    \TabStaff
    \remove Time_signature_engraver

    }
}

<<
  \new TabStaff <<\notes>>
>>
\end{lilypond}
\end{center}

\audioclip{Cierre en \maqam Nahawand.}

\newending{}

\begin{center}
\begin{lilypond}
\include "arabic.ly"
notes = \relative {
mib4:32 re4:32 do16 sol la si do do, do'4
\bar "||"
}

\layout {
  \omit Voice.StringNumber
  \context {
      \TabStaff
      \tabFullNotation
      stringTunings = \stringTuning <do, fa, la, re sol do'>
    }
  indent = #0
  \context {
    \Score
    defaultBarType = ""
    supportNonIntegerFret = ##t
  }
  \context {
    \TabStaff
    \remove Time_signature_engraver

    }
}

<<
  \new TabStaff <<\notes>>
>>
\end{lilypond}
\end{center}

\audioclip{Cierre en \maqam Nahawand.}

\newending{}

\begin{center}
\begin{lilypond}
\include "arabic.ly"
notes = \relative {
sol8 sol8 sol16 lab sib lab sol fa mib re do si do4
\bar "||"
}

\layout {
  \omit Voice.StringNumber
  \context {
      \TabStaff
      \tabFullNotation
      stringTunings = \stringTuning <do, fa, la, re sol do'>
    }
  indent = #0
  \context {
    \Score
    defaultBarType = ""
    supportNonIntegerFret = ##t
  }
  \context {
    \TabStaff
    \remove Time_signature_engraver

    }
}

<<
  \new TabStaff <<\notes>>
>>
\end{lilypond}
\end{center}

\audioclip{Cierre en \maqam Nahawand.}

\newending{}

\begin{center}
\begin{lilypond}
\include "arabic.ly"
notes = \relative {
\time 3/8
do,16 fa' fa fa fa fa do, mib' mib mib mib mib do, re' re re re do, do'4

\bar "||"
}

\layout {
  \omit Voice.StringNumber
  \context {
      \TabStaff
      \tabFullNotation
      stringTunings = \stringTuning <do, fa, la, re sol do'>
    }
  indent = #0
  \context {
    \Score
    defaultBarType = ""
    supportNonIntegerFret = ##t
  }
  \context {
    \TabStaff
    \remove Time_signature_engraver

    }
}

<<
  \new TabStaff <<\notes>>
>>
\end{lilypond}
\end{center}

\audioclip{Cierre en \maqam Nahawand.}

\label{frases}
\section{Frases breves}

Esta sección contiene frases breves que aparecen repetidamente en gran numero de \taqasim, y que pueden considerarse como <<clichés>> dentro de la improvisación con el ud.

Éntiendase esta sección como un simple <<diccionario>> de esas expresiones, de interés para enriquecer el vocabulario a la hora de improvisar.

La mayoría de las frases, al ser breves, usan notas que forman parte de un único \yinsb. Por ello, pueden emplearse siempre que el \maqam en el que se toque contenga dicho \yinsb. Esta es una manera muy útil de aprender las frases, de manera que después puedan reutilizarse en diversos contextos melódicos.

En algunas de las grabaciones correspondientes, se incluye no solo la frase en sí, sino un fragmento más amplio que contiene a esta.

\newphrase{}

Una pequeña frase que aparece con mucha frecuencia, especialmente en las posiciones que permiten el uso de la cuerda al aire.


\begin{center}
\begin{lilypond}
\include "arabic.ly"
notes = \relative {
sol16 (la32 sol) fa16 sol sol
\bar "||"
}

\layout {
  \omit Voice.StringNumber
  \context {
      \TabStaff
      \tabFullNotation
      stringTunings = \stringTuning <do, fa, la, re sol do'>
    }
  indent = #0
  \context {
    \Score
    defaultBarType = ""
    supportNonIntegerFret = ##t
  }
  \context {
    \TabStaff
    \remove Time_signature_engraver

    }
}

<<
  \new TabStaff <<\notes>>
>>
\end{lilypond}
\end{center}

\audioclip{Frase \arabic{phrase}.}

El siguiente ejemplo es un frase ascendente que utiliza el fragmento anterior tanto en la segunda como en la primera cuerda.


\begin{center}
\begin{lilypond}
\include "arabic.ly"
notes = \relative {
mib16 mib fa fa sol (la32 sol) fa16 sol la la sib sib do (re32 do) sib16 do do do do do
\bar "||"
}

\layout {
  \omit Voice.StringNumber
  \context {
      \TabStaff
      \tabFullNotation
      stringTunings = \stringTuning <do, fa, la, re sol do'>
    }
  indent = #0
  \context {
    \Score
    defaultBarType = ""
    supportNonIntegerFret = ##t
  }
  \context {
    \TabStaff
    \remove Time_signature_engraver

    }
}

<<
  \new TabStaff <<\notes>>
>>
\end{lilypond}
\end{center}

\audioclip{Patron utilizando frase \arabic{phrase}.}

\newphrase{}

Otra frase que también aprovecha las cuerdas al aire.


\begin{center}
\begin{lilypond}
\include "arabic.ly"
notes = \relative {
re16 sol fa sol la sib32 la sol16 fa sol sol
\bar "||"
}

\layout {
  \omit Voice.StringNumber
  \context {
      \TabStaff
      \tabFullNotation
      stringTunings = \stringTuning <do, fa, la, re sol do'>
    }
  indent = #0
  \context {
    \Score
    defaultBarType = ""
    supportNonIntegerFret = ##t
  }
  \context {
    \TabStaff
    \remove Time_signature_engraver

    }
}

<<
  \new TabStaff <<\notes>>
>>
\end{lilypond}
\end{center}


\audioclip{Frase \arabic{phrase}.}

\newphrase{}


\begin{center}
\begin{lilypond}
\include "arabic.ly"
notes = \relative {
sol16 sol sol sol la32 sib la sol sol16 sol sol sol
\bar "||"
}

\layout {
  \omit Voice.StringNumber
  \context {
      \TabStaff
      \tabFullNotation
      stringTunings = \stringTuning <do, fa, la, re sol do'>
    }
  indent = #0
  \context {
    \Score
    defaultBarType = ""
    supportNonIntegerFret = ##t
  }
  \context {
    \TabStaff
    \remove Time_signature_engraver

    }
}

<<
  \new TabStaff <<\notes>>
>>
\end{lilypond}
\end{center}

\audioclip{Frase \arabic{phrase}.}

Notese que en estas tres frases se enfatiza el uso de la nota que se da con la cuerda al aire, y que en los tres ejemplos se trata de la nota sol. Como ya dijimos, esta nota es de particular importancia, ya que es el \emph{ghammaz} de buena parte de los \maqamaat habituales.

Jugando con las tres frases anteriores, un inicio y un cierre, y añadiendo algunas otras notas para conectar todo ello, podemos ya hacer frases más largas como la siguiente en \maqam Bayati.

***

\audioclip{Ejemplo de combinación de pequeñas frases.}

\newphrase{}


\begin{center}
\begin{lilypond}
\include "arabic.ly"
notes = \relative {
re'16 mib do re si do lab si sol4
\bar "||"
}

\layout {
  \omit Voice.StringNumber
  \context {
      \TabStaff
      \tabFullNotation
      stringTunings = \stringTuning <do, fa, la, re sol do'>
    }
  indent = #0
  \context {
    \Score
    defaultBarType = ""
    supportNonIntegerFret = ##t
  }
  \context {
    \TabStaff
    \remove Time_signature_engraver

    }
}

<<
  \new TabStaff <<\notes>>
>>
\end{lilypond}
\end{center}


\audioclip{Frase \arabic{phrase}.}


\begin{center}
\begin{lilypond}
\include "arabic.ly"
notes = \relative {
sol8
la16 sib la32 (sib) sol16
la16 sib la32 (sib) sol16
la8 la
\bar "||"
}

\layout {
  \omit Voice.StringNumber
  \context {
      \TabStaff
      \tabFullNotation
      stringTunings = \stringTuning <do, fa, la, re sol do'>
    }
  indent = #0
  \context {
    \Score
    defaultBarType = ""
    supportNonIntegerFret = ##t
  }
  \context {
    \TabStaff
    \remove Time_signature_engraver

    }
}

<<
  \new TabStaff <<\notes>>
>>
\end{lilypond}
\end{center}


\audioclip{Frase \arabic{phrase}.}

\newphrase{}


\begin{center}
\begin{lilypond}
\include "arabic.ly"
notes = \relative {
\time 3/8
sol16 sol16:64 sol16 sol fa sol16
\bar "||"
}

\layout {
  \omit Voice.StringNumber
  \context {
      \TabStaff
      \tabFullNotation
      stringTunings = \stringTuning <do, fa, la, re sol do'>
    }
  indent = #0
  \context {
    \Score
    defaultBarType = ""
    supportNonIntegerFret = ##t
  }
  \context {
    \TabStaff
    \remove Time_signature_engraver

    }
}

<<
  \new TabStaff <<\notes>>
>>
\end{lilypond}
\end{center}


\audioclip{Frase \arabic{phrase}.}

\newphrase{}


\begin{center}
\begin{lilypond}
\include "arabic.ly"
notes = \relative {
fa16 fa16:64 fa16 mib fa sol lab8
\bar "||"
}

\layout {
  \omit Voice.StringNumber
  \context {
      \TabStaff
      \tabFullNotation
      stringTunings = \stringTuning <do, fa, la, re sol do'>
    }
  indent = #0
  \context {
    \Score
    defaultBarType = ""
    supportNonIntegerFret = ##t
  }
  \context {
    \TabStaff
    \remove Time_signature_engraver

    }
}

<<
  \new TabStaff <<\notes>>
>>
\end{lilypond}
\end{center}


\audioclip{Frase \arabic{phrase}.}

\newphrase{}


\begin{center}
\begin{lilypond}
\include "arabic.ly"
notes = \relative {
la32 (sib) do16 sib sib la la sol sol fa sol
\bar "||"
}

\layout {
  \omit Voice.StringNumber
  \context {
      \TabStaff
      \tabFullNotation
      stringTunings = \stringTuning <do, fa, la, re sol do'>
    }
  indent = #0
  \context {
    \Score
    defaultBarType = ""
    supportNonIntegerFret = ##t
  }
  \context {
    \TabStaff
    \remove Time_signature_engraver

    }
}

<<
  \new TabStaff <<\notes>>
>>
\end{lilypond}
\end{center}


\audioclip{Frase \arabic{phrase}.}

\newphrase{}


\begin{center}
\begin{lilypond}
\include "arabic.ly"
notes = \relative {
sib32 lab sol16 fa fa mib mib re re
\bar "||"
}

\layout {
  \omit Voice.StringNumber
  \context {
      \TabStaff
      \tabFullNotation
      stringTunings = \stringTuning <do, fa, la, re sol do'>
    }
  indent = #0
  \context {
    \Score
    defaultBarType = ""
    supportNonIntegerFret = ##t
  }
  \context {
    \TabStaff
    \remove Time_signature_engraver

    }
}

<<
  \new TabStaff <<\notes>>
>>
\end{lilypond}
\end{center}


\audioclip{Frase \arabic{phrase}.}

\newphrase{}


\begin{center}
\begin{lilypond}
\include "arabic.ly"
notes = \relative {
r8 \grace re'32 \glissando mib16 \grace re32 \glissando mib16
do sol do re do sib do4
\bar "||"
}

\layout {
  \omit Voice.StringNumber
  \context {
      \TabStaff
      \tabFullNotation
      stringTunings = \stringTuning <do, fa, la, re sol do'>
    }
  indent = #0
  \context {
    \Score
    defaultBarType = ""
    supportNonIntegerFret = ##t
  }
  \context {
    \TabStaff
    \remove Time_signature_engraver

    }
}

<<
  \new TabStaff <<\notes>>
>>
\end{lilypond}
\end{center}


\audioclip{Frase \arabic{phrase}.}

\newphrase{}


\begin{center}
\begin{lilypond}
\include "arabic.ly"
notes = \relative {
r16 sol16 lab si do re (do) si do4 (do16) sib16 lab sol lab16. sol32 fa8
\bar "||"
}

\layout {
  \omit Voice.StringNumber
  \context {
      \TabStaff
      \tabFullNotation
      stringTunings = \stringTuning <do, fa, la, re sol do'>
    }
  indent = #0
  \context {
    \Score
    defaultBarType = ""
    supportNonIntegerFret = ##t
  }
  \context {
    \TabStaff
    \remove Time_signature_engraver

    }
}

<<
  \new TabStaff <<\notes>>
>>
\end{lilypond}
\end{center}


\audioclip{Frase \arabic{phrase}.}


\newphrase{}


\begin{center}
\begin{lilypond}
\include "arabic.ly"
notes = \relative {
sol8 lab si32 lab si lab si lab si lab lab4 sol
\bar "||"
}

\layout {
  \omit Voice.StringNumber
  \context {
      \TabStaff
      \tabFullNotation
      stringTunings = \stringTuning <do, fa, la, re sol do'>
    }
  indent = #0
  \context {
    \Score
    defaultBarType = ""
    supportNonIntegerFret = ##t
  }
  \context {
    \TabStaff
    \remove Time_signature_engraver

    }
}

<<
  \new TabStaff <<\notes>>
>>
\end{lilypond}
\end{center}


\audioclip{Frase \arabic{phrase}.}


\newphrase{}
%https://www.youtube.com/watch?v=qDXiEb6bwP8

\begin{center}
\begin{lilypond}
\include "arabic.ly"
notes = \relative {
r8 	\grace{re'32} do16 \grace{re32} do16
sib sib	 sol lab4
\bar "||"
}

\layout {
  \omit Voice.StringNumber
  \context {
      \TabStaff
      \tabFullNotation
      stringTunings = \stringTuning <do, fa, la, re sol do'>
    }
  indent = #0
  \context {
    \Score
    defaultBarType = ""
    supportNonIntegerFret = ##t
  }
  \context {
    \TabStaff
    \remove Time_signature_engraver

    }
}

<<
  \new TabStaff <<\notes>>
>>
\end{lilypond}
\end{center}


\newphrase{}

\begin{center}
\begin{lilypond}
\include "arabic.ly"
notes = \relative {
sol16. sol32 sol16 sol16 la32 sib sol la sol16 sol
\bar "||"
}

\layout {
  \omit Voice.StringNumber
  \context {
      \TabStaff
      \tabFullNotation
      stringTunings = \stringTuning <do, fa, la, re sol do'>
    }
  indent = #0
  \context {
    \Score
    defaultBarType = ""
    supportNonIntegerFret = ##t
  }
  \context {
    \TabStaff
    \remove Time_signature_engraver

    }
}

<<
  \new TabStaff <<\notes>>
>>
\end{lilypond}
\end{center}


\audioclip{Frase \arabic{phrase}.}


\label{secuenciasmelodicasvocabulario}
\section{Secuencias melódicas y frases}

Las secuencias melódicas de esta sección no se asocian, por lo general, a un \maqam determinado, y pueden utilizarse en cualquiera de ellos, sin más que adaptar los intervalos y las notas correspondientes. Tambien pueden desplazarse a notas de inicio distintas a las empleadas en la transcripción, manteniéndose dentro del mismo \maqam.

Todas ellas pueden aparecer como parte del desarrollo de una frase más larga, y hacerlo con distinta longitud (el patrón melódico se repite más o menos veces arrancando de distintas notas sucesivas). También pueden emplearse como cierres si se añade un remate al final de la secuencia y se hace terminar esta sobre la tónica del \maqamb. Esto tiene especial sentido en el caso de las secuencias descendentes, pero también las ascendentes pueden emplearse de esta manera.


\audioclip{Secuencia \arabic{sequence}.}

\newsequence{}

\begin{center}
\begin{lilypond}
\include "arabic.ly"
notes = \relative {
do'16. sib32 lab16
sib16. lab32 sol16
lab16. sol32 fa16
sol16. fa32 mib16
\bar "||"
}

\layout {
  \omit Voice.StringNumber
  \context {
      \TabStaff
      \tabFullNotation
      stringTunings = \stringTuning <do, fa, la, re sol do'>
    }
  indent = #0
  \context {
    \Score
    defaultBarType = ""
    supportNonIntegerFret = ##t
  }
  \context {
    \TabStaff
    \remove Time_signature_engraver

    }
}

<<
  \new TabStaff <<\notes>>
>>
\end{lilypond}
\end{center}


\newsequence{}

\begin{center}
\begin{lilypond}
\include "arabic.ly"
notes = \relative {
r16 re'16 sib do la sib sol la fa
do' la sib sol la fa sol misb
sib' sol la fa sol misb fa re8
\bar "||"
}

\layout {
  \omit Voice.StringNumber
  \context {
      \TabStaff
      \tabFullNotation
      stringTunings = \stringTuning <do, fa, la, re sol do'>
    }
  indent = #0
  \context {
    \Score
    defaultBarType = ""
    supportNonIntegerFret = ##t
  }
  \context {
    \TabStaff
    \remove Time_signature_engraver

    }
}

<<
  \new TabStaff <<\notes>>
>>
\end{lilypond}
\end{center}


\audioclip{Secuencia \arabic{sequence}.}

\newsequence{}

\begin{center}
\begin{lilypond}
\include "arabic.ly"
notes = \relative {
do'16 do la32 la la16
sisb sisb sol32 sol sol16
la la fa32 fa fa16
sol sol misb32 misb misb16
\bar "||"
}

\layout {
  \omit Voice.StringNumber
  \context {
      \TabStaff
      \tabFullNotation
      stringTunings = \stringTuning <do, fa, la, re sol do'>
    }
  indent = #0
  \context {
    \Score
    defaultBarType = ""
    supportNonIntegerFret = ##t
  }
  \context {
    \TabStaff
    \remove Time_signature_engraver

    }
}

<<
  \new TabStaff <<\notes>>
>>
\end{lilypond}
\end{center}


\audioclip{Secuencia \arabic{sequence}.}

\newsequence{}

\begin{center}
\begin{lilypond}
\include "arabic.ly"
notes = \relative {
sib16.:128 la32 sol16 sib la sol
la16.:128 sol32 fa16 la sol fa
sol16.:128 fa32 misb16 sol fa misb
\bar "||"
}

\layout {
  \omit Voice.StringNumber
  \context {
      \TabStaff
      \tabFullNotation
      stringTunings = \stringTuning <do, fa, la, re sol do'>
    }
  indent = #0
  \context {
    \Score
    defaultBarType = ""
    supportNonIntegerFret = ##t
  }
  \context {
    \TabStaff
    \remove Time_signature_engraver

    }
}

<<
  \new TabStaff <<\notes>>
>>
\end{lilypond}
\end{center}


\audioclip{Secuencia \arabic{sequence}.}

\newsequence{}

\begin{center}
\begin{lilypond}
\include "arabic.ly"
notes = \relative {
\tuplet 3/2 {re8 do do}
\tuplet 3/2 {mib re re}
\tuplet 3/2 {fa mib mib}
\tuplet 3/2 {sol fa fa}
\tuplet 3/2 {la sol sol}
\bar "||"
}

\layout {
  \omit Voice.StringNumber
  \context {
      \TabStaff
      \tabFullNotation
      stringTunings = \stringTuning <do, fa, la, re sol do'>
    }
  indent = #0
  \context {
    \Score
    defaultBarType = ""
    supportNonIntegerFret = ##t
  }
  \context {
    \TabStaff
    \remove Time_signature_engraver

    }
}

<<
  \new TabStaff <<\notes>>
>>
\end{lilypond}
\end{center}


\audioclip{Secuencia \arabic{sequence}.}

\newsequence{}

\begin{center}
\begin{lilypond}
\include "arabic.ly"
notes = \relative {
do16 re mib mib32 mib mib16
re mib fa fa32 fa fa16
mib fa sol sol32 sol sol16
fa sol lab lab32 lab lab16
\bar "||"
}

\layout {
  \omit Voice.StringNumber
  \context {
      \TabStaff
      \tabFullNotation
      stringTunings = \stringTuning <do, fa, la, re sol do'>
    }
  indent = #0
  \context {
    \Score
    defaultBarType = ""
    supportNonIntegerFret = ##t
  }
  \context {
    \TabStaff
    \remove Time_signature_engraver

    }
}

<<
  \new TabStaff <<\notes>>
>>
\end{lilypond}
\end{center}


\audioclip{Secuencia \arabic{sequence}.}


\newsequence{}

\begin{center}
\begin{lilypond}
\include "arabic.ly"
notes = \relative {
lab16 lab sol sol fa fa mib mib
sol sol fa fa mib mib re re
fa fa mib mib re re do do
mib mib re re do do sib sib
\bar "||"
}

\layout {
  \omit Voice.StringNumber
  \context {
      \TabStaff
      \tabFullNotation
      stringTunings = \stringTuning <do, fa, la, re sol do'>
    }
  indent = #0
  \context {
    \Score
    defaultBarType = ""
    supportNonIntegerFret = ##t
  }
  \context {
    \TabStaff
    \remove Time_signature_engraver

    }
}

<<
  \new TabStaff <<\notes>>
>>
\end{lilypond}
\end{center}


\audioclip{Secuencia \arabic{sequence}.}

\newsequence{}

\begin{center}
\begin{lilypond}
\include "arabic.ly"
notes = \relative {
mib'16-> re do do
re-> do sib sib
do-> sib lab lab
sib-> lab sol sol
\bar "||"
}

\layout {
  \omit Voice.StringNumber
  \context {
      \TabStaff
      \tabFullNotation
      stringTunings = \stringTuning <do, fa, la, re sol do'>
    }
  indent = #0
  \context {
    \Score
    defaultBarType = ""
    supportNonIntegerFret = ##t
  }
  \context {
    \TabStaff
    \remove Time_signature_engraver

    }
}

<<
  \new TabStaff <<\notes>>
>>
\end{lilypond}
\end{center}


\audioclip{Secuencia \arabic{sequence}.}

\newsequence{}

\begin{center}
\begin{lilypond}
\include "arabic.ly"
notes = \relative {
do'16 sib lab lab sol sol
sib lab sol sol fa fa
lab sol fa fa mib mib
sol fa mib mib re re
\bar "||"
}

\layout {
  \omit Voice.StringNumber
  \context {
      \TabStaff
      \tabFullNotation
      stringTunings = \stringTuning <do, fa, la, re sol do'>
    }
  indent = #0
  \context {
    \Score
    defaultBarType = ""
    supportNonIntegerFret = ##t
  }
  \context {
    \TabStaff
    \remove Time_signature_engraver

    }
}

<<
  \new TabStaff <<\notes>>
>>
\end{lilypond}
\end{center}


\audioclip{Secuencia \arabic{sequence}.}

\newsequence{}

\begin{center}
\begin{lilypond}
\include "arabic.ly"
notes = \relative {
\tuplet 3/2{fa16 sol lab} sol8
\tuplet 3/2{mib16 fa sol} fa8
\tuplet 3/2{re16 mib fa} mib8
\tuplet 3/2{do16 re mib} re8
\bar "||"
}

\layout {
  \omit Voice.StringNumber
  \context {
      \TabStaff
      \tabFullNotation
      stringTunings = \stringTuning <do, fa, la, re sol do'>
    }
  indent = #0
  \context {
    \Score
    defaultBarType = ""
    supportNonIntegerFret = ##t
  }
  \context {
    \TabStaff
    \remove Time_signature_engraver

    }
}

<<
  \new TabStaff <<\notes>>
>>
\end{lilypond}
\end{center}


\audioclip{Secuencia \arabic{sequence}.}

\newsequence{}

\begin{center}
\begin{lilypond}
\include "arabic.ly"
notes = \relative {
sib8 \tuplet 3/2 {sib16 lab sol}
lab8 \tuplet 3/2 {lab16 sol fa}
sol8 \tuplet 3/2 {sol16 fa mib}
fa8 \tuplet 3/2 {fa16 mib re} mib8
\bar "||"
}

\layout {
  \omit Voice.StringNumber
  \context {
      \TabStaff
      \tabFullNotation
      stringTunings = \stringTuning <do, fa, la, re sol do'>
    }
  indent = #0
  \context {
    \Score
    defaultBarType = ""
    supportNonIntegerFret = ##t
  }
  \context {
    \TabStaff
    \remove Time_signature_engraver

    }
}

<<
  \new TabStaff <<\notes>>
>>
\end{lilypond}
\end{center}


\audioclip{Secuencia \arabic{sequence}.}

\newsequence{}

\begin{center}
\begin{lilypond}
\include "arabic.ly"
notes = \relative {
lab8 sol16 fa fa16.mib32 re8
sol8 fa16 mib mib16.re32 do8
fa8 mib16 re re16.do32 si8
\bar "||"
}

\layout {
  \omit Voice.StringNumber
  \context {
      \TabStaff
      \tabFullNotation
      stringTunings = \stringTuning <do, fa, la, re sol do'>
    }
  indent = #0
  \context {
    \Score
    defaultBarType = ""
    supportNonIntegerFret = ##t
  }
  \context {
    \TabStaff
    \remove Time_signature_engraver

    }
}

<<
  \new TabStaff <<\notes>>
>>
\end{lilypond}
\end{center}


\audioclip{Secuencia \arabic{sequence}.}

\newsequence{}

\begin{center}
\begin{lilypond}
\include "arabic.ly"
notes = \relative {
do'16 sisb16. la32 sol16
sisb la16. sol32 fa16
la sol16. fa32 misb16
sol fa16. misb32 re16
\bar "||"
}

\layout {
  \omit Voice.StringNumber
  \context {
      \TabStaff
      \tabFullNotation
      stringTunings = \stringTuning <do, fa, la, re sol do'>
    }
  indent = #0
  \context {
    \Score
    defaultBarType = ""
    supportNonIntegerFret = ##t
  }
  \context {
    \TabStaff
    \remove Time_signature_engraver

    }
}

<<
  \new TabStaff <<\notes>>
>>
\end{lilypond}
\end{center}


\audioclip{Secuencia \arabic{sequence}.}
